% ============================================================
% T3：Hitchin 模空间与超凯勒几何（REF-TERMS group C，15 条）
% 依赖 rosetta/preamble.tex：\track、word 环境、\mathdef 等七个区块
% ============================================================

\track{T3\quad Hitchin 模空间与超凯勒几何}{本轨道 15 条：Hitchin 模空间是两个世界的交汇处——它既是 4 维规范理论的低能舞台，又是几何 Langlands 的几何载体。}

\begin{word}{Hitchin 模空间}{Hitchin moduli space}{hitchin-moduli}{§4.1 p.45；§4 p.38--55}
\mathdef{设 $G$ 为紧半单 Lie 群，$C$ 为亏格 $g$ 的紧黎曼面，$E\to C$ 为主 $G$-丛，$A$ 为 $E$ 上联络，$\varphi\in\Omega^1(C,\mathrm{ad}E)$ 为伴随值 1-形式。Hitchin 方程是 $F_A-\varphi\wedge\varphi=0$、$d_A\varphi=0$、$d_A^*\varphi=0$，其中 $F_A$ 为曲率、$d_A$ 为协变外微分、$d_A^*$ 为其形式伴随（把 $\varphi$ 换成 $i\varphi$ 即得常见的 $F_A+\varphi\wedge\varphi=0$ 写法；原文式 (4.1) 用前一种）。$\mathcal M_H(G,C)$ 定义为这些方程的解模去规范变换所得的模空间。当 $g>1$ 时它是复流形，$\dim_{\mathbb C}\mathcal M_H=(2g-2)\dim G$，即 $2\dim_{\mathbb C}\mathcal M_H=4(g-1)\dim G$。}
\mathex{$G=SU(2)$、$C=\mathbb P^1$（$g=0$）：$H^0(\mathbb P^1,K_C)=0$ 迫使 $\varphi=0$，方程化为 $F_A=0$；而 $\mathbb P^1$ 单连通，只有平凡平坦联络，故 $\mathcal M_H$ 退化为一个点（粗模空间意义下）。$g\ge2$ 才有丰富几何：$G=SU(2)$、$g=2$ 时 $\dim_{\mathbb C}\mathcal M_H=6$，Hitchin 纤维化的底空间 $B=H^0(C,K_C^2)$ 复维数为 $3g-3=3$。}
\physdef{把四维扭化 $\mathcal N=4$ 超对称杨--米尔斯理论紧化到 $M=\Sigma\times C$，并令 $C$ 远小于 $\Sigma$：$C$ 上的高阶模式被积掉，$\Sigma$ 上的低能有效理论是靶空间为 $\mathcal M_H$ 的二维 sigma 模型。$\mathcal M_H$ 参数化四维拓扑扭化理论在 $C$ 上的 BPS 位形，是 $C$ 上 Hitchin 方程解的模空间。}
\physex{$G=SU(2)$ 时 $\mathcal M_H$ 是二维 sigma 模型的靶空间，其三个复结构让同一个二维理论呈现为不同拓扑场论。$G=U(1)$ 或 $U(N)$ 时每个解还带未破缺的 $U(1)$ 对称性，低能描述是 sigma 模型与超对称 $U(1)$ 规范场论的乘积。}
\bridge{数学侧的 $\mathcal M_H$ 与物理侧的低能靶空间是同一个实流形：它的点在复结构 $I$ 下是半稳定 Higgs 丛、在复结构 $J$ 下是平坦 $G_{\mathbb C}$-丛；物理上每个点代表二维场论的一个真空（一个 BPS 位形）。$\mathcal M_H$ 的奇点对应有连续未破缺规范对称性的可约解，也正是低能描述失效处（§4.1、§4.2）。}
\intuit{把 $\mathcal M_H$ 想成「$C$ 上所有 Hitchin 解收进的空间」：每个点是一种解，二维低能场论就是这个空间里的 sigma 模型。直觉靠得住，是因为 Hitchin 方程是椭圆方程组，解空间有限维且自带超凯勒结构。}
\pitfall{$\mathcal M_H$ 不是 $G$-丛模空间 $\mathcal M$：后者只是 $\phi=0$ 的子空间。也不要把它与四维理论的真空模空间或「Hitchin 系统的相空间」混为一谈——后者指的是 $\mathcal M_H$ 配上辛形式后的可积系统，而「相空间」一词还暗示实哈密顿系统，需要另选实切片。}
\end{word}

\begin{word}{Higgs 丛}{Higgs bundle}{higgs-bundle}{§4.2 p.49}
\mathdef{Higgs 丛是 $C$ 上的对 $(E,\varphi)$：$E$ 是 $C$ 上的全纯 $G_{\mathbb C}$-主丛（$G_{\mathbb C}$ 为 $G$ 的复化），$\varphi$ 是取值于伴随丛 $\mathrm{ad}E$、系数在典范线丛 $K_C$ 的全纯 1-形式，即 $\varphi\in H^0(C,K_C\otimes\mathrm{ad}E)$；等价地，$\varphi$ 是 $\mathrm{ad}E$-值 $(1,0)$-形式且 $\bar\partial_E\varphi=0$。取 $G=SU(2)$ 时 $E$ 是秩 2 全纯向量丛，$\varphi:E\to E\otimes K_C$ 是 $\mathcal O_C$-线性映射（在 $\mathfrak{sl}_2$ 中取值）。}
\mathex{$G=SU(2)$、$C=\mathbb P^1$：$K_C=\mathcal O(-2)$ 没有非零全纯截面，故必有 $\varphi=0$，Higgs 丛退化为普通全纯丛；$\mathbb P^1$ 上度为 $0$ 的半稳定 $SU(2)$ 丛只有平凡丛 $\mathcal O\oplus\mathcal O$。$g\ge2$ 时情形丰富：固定一个稳定丛 $E$，全部可能的 $\varphi$ 构成复维数 $H^0(C,K_C\otimes\mathrm{ad}E)$ 的线性空间，$SU(2)$ 情形维数是 $3(g-1)$。}
\physdef{Higgs 丛是 $\mathcal M_H$ 在复结构 $I$ 下作为复流形的「点」：$A_{\bar z}$ 决定全纯结构 $E$，$\phi$ 的 $(1,0)$ 分量 $\phi_z\,dz$ 决定 $\varphi$。物理上它把一个二维低能真空（一个 Hitchin 方程解）按复结构 $I$ 解码成全纯数据，是四维扭化理论 BPS 方程在 $C$ 上的全纯描述。}
\physex{$G=SU(2)$、$C$ 亏格 2：每个点是一个半稳定对 $(E,\varphi)$，其中 $\varphi$ 有 $3$ 个复自由度。$\mathcal M_H$ 的复维数 $6=(2g-2)\dim G$ 恰好可读成「丛的模 $3$ 维 + Higgs 场的模 $3$ 维」。}
\bridge{$\mathcal M_H$ 在复结构 $I$ 下恰是半稳定 Higgs 丛的模空间（含严格半稳定的某些等价类），即原文 §4.2 的操作 (II)$'$；忘掉 $\varphi$ 得到「第二纤维化」到稳定丛模空间 $\mathcal M$，而 $T^*\mathcal M$ 是 $\mathcal M_H$ 的 Zariski 开集（§4.3）。}
\intuit{把 Higgs 丛想成「丛 + 一个全纯的伴随矩阵值 1-形式」：丛负责规范场，$\varphi$ 负责 Higgs 场在全纯世界的部分。直觉靠得住，是因为在复结构 $I$ 中这两组数据恰好是 $\mathcal M_H$ 的全纯坐标。}
\pitfall{不要把 Higgs 丛等同于「Higgs 场的真空模空间」：真空模空间（靶空间）是 $\mathcal M_H$ 整个实流形，Higgs 丛只是它在复结构 $I$ 下的描述。也不要把 $\varphi$ 当作普通截面：它带 $K_C$ 的系数，换坐标时会按微分形式变换，$SU(2)$ 情形下它是 $E\to E\otimes K_C$ 的映射而不是 $E\to E$ 的自同态。}
\end{word}

\begin{word}{Higgs 场}{Higgs field}{higgs-field}{§4.2 p.49}
\mathdef{Higgs 场就是 Higgs 丛 $(E,\varphi)$ 中的 $\varphi$：$C$ 上取值于 $\mathrm{ad}E$、系数在 $K_C$ 的全纯 1-形式，$\varphi\in H^0(C,K_C\otimes\mathrm{ad}E)$。局部写 $\varphi=f(z)\,dz$，其中 $f$ 取值于 $\mathrm{ad}E=\mathrm{End}_0E$；换坐标卡时 $f$ 按 $\mathrm{ad}E$ 的转移函数与 $dz$ 的变换同时变化，故 $\varphi$ 的整体性由 $K_C\otimes\mathrm{ad}E$ 保证。$G=SU(N)$ 且 $E$ 平凡时，$\varphi$ 就是取值于 $\mathfrak{sl}(N,\mathbb C)$ 的全纯 1-形式矩阵。}
\mathex{$G=SU(2)$、$g\ge2$：固定稳定丛 $E$ 后，Higgs 场取值于 $H^0(C,K_C\otimes\mathrm{ad}E)$，复维数 $3(g-1)$；$g=2$ 时是 3 维，正好是 $\mathcal M_H$ 复维数的一半。$G=U(1)$ 时 $\varphi\in H^0(C,K_C)$ 是普通全纯 1-形式，复维数 $g$，此时 Higgs 丛模空间是 $T^*\mathrm{Jac}(C)$。}
\physdef{四维扭化 $\mathcal N=4$ 理论中，除规范场 $A$ 外还有伴随值 1-形式场 $\phi$（以及复标量 $\sigma$）。在 $M=\Sigma\times C$ 上，$\phi$ 沿 $C$ 的分量在 $\Sigma$ 上表现为标量；其 $(1,0)$ 部分 $\varphi=\phi_z\,dz$ 就是 Higgs 场，其余分量在 Hitchin 方程中被解掉，而 $\sigma$ 在 Hitchin 解中取零。}
\physex{$G=SU(2)$、$C$ 亏格 2：$\varphi$ 是一个 $2\times2$ 无迹矩阵值的全纯 1-形式，在 $C$ 上有 3 个复自由度。二维 sigma 模型看它时，$\varphi$ 的 3 个复分量给出靶空间 $\mathcal M_H$ 的 12 个实坐标中的 6 个。}
\bridge{物理侧的 $\phi$ 与数学侧的 $\varphi$ 是同一个场在两个复结构下的两种读法：在复结构 $I$ 中读 $A_{\bar z}$ 与 $\phi_z$，$\phi_z\,dz$ 即 Higgs 场；在复结构 $J$ 中读 $\mathcal A=A+i\phi$，$\phi$ 被并进复平坦联络（§4.1--§4.2）。}
\intuit{把 Higgs 场想成「丛上的一个矩阵值微分形式」：它在每一点给出一个无迹矩阵，沿 $C$ 全纯地变化。直觉的边界是：它不是标量场，而是 1-形式，所以它的「值」依赖切向量。}
\pitfall{最常见的错：把这里的 Higgs 场与四维理论的复标量 $\sigma$ 或与 Higgs 机制里的标量混为一谈。$\varphi$ 是沿 $C$ 的 $(1,0)$-形式；在 Hitchin 解中 $\sigma=0$、$\varphi\neq0$。另一个错是把 $\varphi$ 当作四维时空中的场：它是紧化后 $C$ 上的数据，在二维理论里才表现为标量。}
\end{word}

\begin{word}{Hitchin 对}{Hitchin pair}{hitchin-pair}{§4.2 p.49}
\mathdef{Hitchin 对是 $C$ 上的对 $(E,\varphi)$：$E$ 为全纯 $G_{\mathbb C}$-丛，$\varphi\in H^0(C,K_C\otimes\mathrm{ad}E)$，且满足全纯方程 $\bar D\varphi=0$（$\bar D$ 是 $E$ 上联络的 $(0,1)$ 部分，即 $\bar\partial_E$）。在 $\mathcal M_H$ 的构造中，$E$ 由 $A_{\bar z}$ 决定，$\varphi=\phi_z\,dz$，于是 $\bar D\varphi=0$ 正是 Hitchin 方程中的复方程。要得到模空间，还需对 $(E,\varphi)$ 加上稳定性条件（见「稳定性」条）并取适当的等价类。}
\mathex{$G=SU(2)$、$g\ge2$：Hitchin 对是半稳定 Higgs 丛的候选者，每个稳定的对给出 $\mathcal M_H$ 的一点（某些情形下是轨形点）。$G=SU(2)$、$C=\mathbb P^1$：唯一可能的对是平凡丛配 $\varphi=0$。$g=2$、$G=SU(2)$：$\varphi=0$ 的对落在零截面 $T^*\mathcal M$ 上，$\varphi\neq0$ 的对才是 $\mathcal M_H$ 里新出现的点。}
\physdef{物理上 Hitchin 对是「把 Hitchin 方程的解按复结构 $I$ 记录的半成品」：先只解全纯方程 $\bar D\varphi=0$，暂不管 moment map 条件。四维扭化理论的 BPS 位形在复结构 $I$ 下给出的正是这种全纯数据（原文 §4.2 的操作 (II)）。}
\physex{$G=SU(2)$、$C$ 亏格 2：Hitchin 对的空间在稳定对处与 $\mathcal M_H$ 一致；若把丛固定而让 $\varphi$ 自由变化，得到的是 $T^*\mathcal M$ 在其上的纤维，复维数 $3(g-1)=3$。}
\bridge{$I$ 结构下 $\mathcal M_H$ 等同于稳定 Hitchin 对的模空间（Hitchin 1987，原文 §4.2）：全纯 moment map $\nu_I=\mu_J+i\mu_K$ 的零点给 $\bar D\varphi=0$，再置 $\mu_I=0$ 并加稳定性，就与超凯勒商的操作 (I) 一致。}
\intuit{把 Hitchin 对想成「半成品解」：全纯方程容易满足、也容易用代数几何理解；剩下的实方程（moment map）与稳定性只是把它校正成真正的解。靠得住的原因是 Kempf--Ness 型论证：几乎每条复规范变换轨道恰含一个 moment map 零点轨道。}
\pitfall{Hitchin 对不自动是 $\mathcal M_H$ 的点：不稳定的对必须扔掉，严格半稳定的对还要作额外的等价。也不要把「满足全纯方程的对」直接当成模空间——操作 (II) 只给近似，加上稳定性后的操作 (II)$'$ 才与 $\mathcal M_H$ 一致。}
\end{word}

\begin{word}{稳定性}{stability}{stability}{§4.2 p.50}
\mathdef{$G=SU(2)$ 情形：若线丛 $\mathcal L\subset E$ 满足 $\varphi(\mathcal L)\subset\mathcal L\otimes K_C$，称 $\mathcal L$ 为 $\varphi$-不变。若每个 $\varphi$-不变线丛的 $\deg\mathcal L<0$，称 $(E,\varphi)$ 稳定；若 $\deg\mathcal L\le0$，称半稳定（此处 $E$ 秩 2、$\deg E=0$）。一般 $G$：把 $E$ 的结构群约化到极大抛物子群 $P$ 的 $\varphi$-不变约化，用 $P$ 中特殊 $U(1)$ 给出的第一 Chern 类，要求它对每个这样的约化都为负（稳定）或非正（半稳定）。等价的斜率写法：每个 $\varphi$-不变子丛 $F\subset E$ 满足 $\mu(F)=\deg F/\mathrm{rk}F<\mu(E)$，半稳定用 $\le$。}
\mathex{$G=SU(2)$、$C$ 亏格 2：稳定 Hitchin 对构成 $\mathcal M_H$ 的开子集，不稳定对余维很高；稳定对给出光滑点，严格半稳定的对给出更坏的奇点。$\varphi=0$ 时条件退化为通常的丛稳定性，故稳定丛模空间 $\mathcal M$ 嵌入 $\mathcal M_H$。}
\physdef{物理上稳定性是「可实现性条件」：一个 Higgs 丛能由 Hitchin 方程的解代表，当且仅当它是（多）稳定的（Kobayashi--Hitchin 对应）。不稳定对在二维低能 sigma 模型的靶空间里没有对应的真空；它是代数几何里被扔掉的数据。}
\physex{$G=SU(2)$、$C$ 亏格 2：$\mathcal M_H$ 的每一点都对应稳定（或严格半稳定）的 Higgs 丛；奇点处正是严格半稳定对导致稳定子变大，低能描述在那些点失效（§4.2）。}
\bridge{$\mathcal M_H$ 只参数化稳定 Hitchin 对与严格半稳定对的某些等价类（原文 §4.2）：这正是操作 (II)$'$ 与操作 (I) 一致的原因，也解释了为什么 $\mathcal M_H$ 既有光滑点又有奇点。}
\intuit{稳定性像「防退化保险」：若存在斜率超标的 $\varphi$-不变子对象，Hitchin 方程就无解；排除它们，商才分离、有限维、几何良好。直觉的边界是：条件必须对所有 $\varphi$-不变对象验证，而不只是子丛。}
\pitfall{两个易错点：其一，方向容易记反——稳定要求第一 Chern 类为负（半稳定非正），不是正；其二，只检查子丛不够，一般 $G$ 要考虑抛物约化。忘掉稳定性只得到拟模空间或叠，不能当作 $\mathcal M_H$。}
\end{word}

\begin{word}{超凯勒流形}{hyper-Kähler manifold}{hyperkahler-manifold}{§4.1 p.42}
\mathdef{超凯勒流形是四元组 $(X,g,I,J,K)$：$X$ 为实流形，$g$ 为黎曼度规，$I,J,K$ 为三个复结构，满足 $I^2=J^2=K^2=IJK=-1$（等价于 $IJ=-JI=K$ 及循环置换），且 $g$ 对每个复结构都是 Kähler 度规，即三个 2-形式 $\omega_I=g(I\cdot,\cdot)$、$\omega_J$、$\omega_K$ 都闭且非退化。由此 $\dim_{\mathbb R}X=4n$、$\dim_{\mathbb C}X=2n$；$\Omega_I=\omega_J+i\omega_K$ 关于 $I$ 是 $(2,0)$ 型闭形式，给出一支全纯辛形式。}
\mathex{最简单的例子是平坦 $\mathbb R^4=\mathbb C^2$：取 $I,J,K$ 为四元数单位 $i,j,k$ 的左乘，三个辛形式都是常系数形式。紧例子是 K3 曲面（复维 2）；非紧例子是 $T^*\mathbb P^1$（Eguchi--Hanson 空间）。本文的主角 $\mathcal M_H(SU(2),C)$ 在 $g=2$ 时是复 6 维（实 12 维）超凯勒流形，其超凯勒结构在 §4.1 由无穷维空间 $\mathcal U$ 上的平坦结构商下来。}
\physdef{超凯勒流形是二维 $(4,4)$ 超对称 sigma 模型的靶空间：三个复结构对应三种超荷与扭化选择，使同一个 sigma 模型在不同复结构下呈现为不同拓扑场论（A-模型与 B-模型）。在 KW06 中，四维扭化理论紧化到 $C$ 后的低能有效理论恰是这类 sigma 模型，靶空间为 $\mathcal M_H$。}
\physex{$G=SU(2)$、$C$ 亏格 2：靶空间实维 12；§5 表明不同 twist 参数 $t$ 给出不同的二维拓扑场论——$t=\pm i$ 给复结构 $J$ 中的 B-模型，$t=\pm1$ 给复结构 $K$ 中的 A-模型，而 $t$ 为实数时永远是 A-模型。}
\bridge{数学侧的超凯勒结构就是物理侧 $(4,4)$ 超对称的几何翻译：三个复结构 $I,J,K$ 对应三个超荷方向，三个辛形式 $\omega_I,\omega_J,\omega_K$ 对应不同扭化下的作用量数据（§4.1、§5）；$I$ 下读出 Higgs 丛，$J$ 下读出平坦丛，$K=IJ$。}
\intuit{超凯勒 = 「一个流形上同时住着三个 Kähler 结构，像四元数 $i,j,k$ 一样互相转动」：任取单位线性组合 $aI+bJ+cK$（$a^2+b^2+c^2=1$）仍是复结构。直觉靠得住，是因为这正是二维 $(4,4)$ 超对称要求靶空间具有 Sp($n$) 和乐。}
\pitfall{只有三个复结构 $I,J,K$ 还不够：必须要求同一个度规 $g$ 对每个复结构都是 Kähler 的（三个 2-形式都闭）。复维必须为偶数（实维是 4 的倍数）。也不要把超凯勒与 hypercomplex（三个复结构但无相容度规）或与普通 Kähler 流形混为一谈。}
\end{word}

\begin{word}{超凯勒商}{hyper-Kähler quotient}{hyperkahler-quotient}{§4.1 p.40；§4.1 p.45}
\mathdef{设 $(X,g,I,J,K)$ 为超凯勒流形，Lie 群 $H$ 以保持 $g,I,J,K$ 的方式作用，并有三个 moment map $\mu_I,\mu_J,\mu_K:X\to\mathfrak h^*$（定义见「moment map」条）。超凯勒商定义为 $X/\!/\!/H:=\vec\mu^{-1}(0)/H$，即三个 moment map 同时为零的解集模 $H$。它是超凯勒流形（Hitchin--Karlhede--Lindström--Roček）：三个 moment map 的零点继承 $X$ 的超凯勒结构，商掉保结构的群后结构仍然存在。}
\mathex{本文的核心例子：$X=\mathcal U$ 是 $C$ 上所有 $(A,\phi)$ 构成的无穷维仿射空间（联络 + 伴随值 1-形式），带平坦超凯勒度规；$H=\mathcal G=\mathrm{Maps}(C,G)$ 是规范群。三个 moment map 恰给出 $F-\phi\wedge\phi=0$ 与 $D\phi=0$、$D\star\phi=0$（式 (4.11)），故 $\mathcal M_H=\mathcal U/\!/\!/\mathcal G$。}
\physdef{在二维 $(4,4)$ 超对称规范理论中，超多重态的参数空间是超凯勒流形 $X$，规范群 $H$ 作用其上；低能下超多重态被约束到三个 D-项方程（moment map 零点），规范多重态获得质量被积掉，故低能有效理论是靶空间为 $X/\!/\!/H$ 的 sigma 模型。}
\physex{把四维扭化理论紧化到 $C$：所得二维理论有 $(4,4)$ 超对称（16 个超荷中保 8 个），$H=\mathrm{Maps}(C,G)$、$X=\mathcal U$；于是低能理论自动是靶空间 $\mathcal U/\!/\!/\mathcal G=\mathcal M_H$ 的 sigma 模型。无穷维的商能给出有限维流形，是因为 Hitchin 方程椭圆。}
\bridge{超凯勒商是「Hitchin 方程解模规范变换」的精确几何表述：$\vec\mu=0$ 就是 Hitchin 方程的三个实方程，除以 $\mathcal G$ 就是模规范变换（原文 §4.1 式 (4.11) 之后即证此点）。物理上它就是二维 $(4,4)$ 理论取低能极限、积掉规范多重态的过程。}
\intuit{超凯勒商 = 普通 Kähler 商的「三倍 D-项」版本：三个 moment map 同时归零再商，把三个方向的冗余一起消掉，商下来的空间对三个复结构都 Kähler。直觉的边界是：这里除的是无穷维群，能不能得到好空间要靠方程椭圆。}
\pitfall{不是普通 Kähler 商 $X/\!/H=\mu^{-1}(0)/H$（只取一个 moment map）：只置 $\mu_I=0$ 得到复几何商（半稳定 Higgs 丛），三个一起归零才是 Hitchin 方程。另一个错是以为商掉有限维群——这里除的是无穷维规范群，关键在 Hitchin 方程椭圆、解空间有限维。}
\end{word}

\begin{word}{复结构三元组}{complex structures $I,J,K$}{complex-structures-IJK}{§4.1 p.42}
\mathdef{$\mathcal M_H$ 上的三个复结构由 $\mathcal U$ 上平移不变的复结构下降得到：$I$ 使 $A_{\bar z}$ 与 $\phi_z$ 成为全纯坐标（$\bar\partial_A$ 定义全纯结构、$\phi_z$ 全纯），$J$ 使 $\mathcal A=A+i\phi$ 成为全纯联络（即复平坦联络），$K=IJ$ 为第三个。三者满足四元数关系 $I^2=J^2=K^2=IJK=-1$，且都与超凯勒度规相容（是超凯勒结构的一部分，不是任意三个复结构）。}
\mathex{$G=SU(2)$：$I$ 下的点是半稳定 Higgs 丛 $(E,\varphi)$；$J$ 下的点是同态 $\vartheta:\pi_1(C)\to SL(2,\mathbb C)$ 的共轭类，即 $Y=\mathrm{Hom}(\pi_1(C),SL(2,\mathbb C))/\mathrm{conj}$（含稳定同态与半稳定等价类）。$J$ 与 $-J$ 只依赖 $G_{\mathbb C}$ 与 $\pi_1(C)$，不依赖 $C$ 的复结构；$I$ 依赖 $C$ 的复结构。}
\physdef{在二维 sigma 模型中三个复结构对应不同拓扑场论：复结构 $J$ 中的 B-模型出现在 twist 参数 $t=\pm i$；复结构 $K$ 中的 A-模型出现在 $t=\pm1$（A-模型真正依赖的是辛形式 $\omega_K$）；复结构 $I$ 中的 B-模型是几何 Langlands 数学侧的出发点，但四维理论给出的族 $t\in\mathbb P^1_g$ 里没有它——该族在实 $t$ 处给 A-模型，$t=0,\infty$ 只给 $I$ 或 $-I$ 中的 A-模型；$I$ 中的 B-模型要靠二维理论的完整参数 $(w_+,w_-)\in\mathbb P^1_h\times\mathbb P^1_h$ 才实现（§5）。}
\physex{$t=\pm i$ 时得到复结构 $J$ 或 $-J$ 中的 B-模型，而 $J$ 不依赖 $C$ 的复结构；实 $t$ 给 A-模型，其中 $t=\pm1$ 落在复结构 $K$ 或 $-K$，$t=0,\infty$ 落在 $I$ 或 $-I$。故沿 $\mathbb P^1_g$ 上的 $t$，二维 TQFT 在 $t=\pm i$ 处为 B-模型、在实 $t$ 处为 A-模型。}
\bridge{$S$-对偶/镜像对称交换的是复结构 $I$ 与 $J$ 两侧（A-模型与 B-模型靶空间互换）：$B$-模型 on $Y({}^LG)$ ↔ $A$-模型 on $\mathcal M_H(G)$（§5）。若把 $I$ 与 $J$ 弄混，膜与 Hitchin 纤维的对应会全部颠倒——译本 §4.1 的译注特别强调这一点。}
\intuit{同一实流形上的三副「眼镜」：戴 $I$ 看是 Higgs 丛，戴 $J$ 看是平坦丛，戴 $K$ 看是第三副。三副眼镜由四元数关系联系，故几何与物理的翻译可以在其间自由切换，而流形本身不变。}
\pitfall{最易错的是把 $I$ 与 $J$ 的物理含义弄反：$J$ 的全纯函数是 $A+i\phi$（复平坦联络），故 $J$ 对应平坦丛与 B-模型；$I$ 的全纯函数是 $A_{\bar z}$ 与 $\phi_z$，故 $I$ 对应 Higgs 丛。另注意 $J$ 不依赖 $C$ 的复结构而 $I$ 依赖；$K$ 不是独立的新结构，而是 $K=IJ$。}
\end{word}

\begin{word}{复结构 $\mathbb P^1_g$}{twistor $\mathbb P^1$}{twistor-p1}{§4.1 p.43}
\mathdef{超凯勒流形 $(X,I,J,K)$ 上有一族由 $\mathbb P^1$ 参数化的复结构：对 $w\in\mathbb C^*$ 令
$I_w=\frac{1-|w|^2}{1+|w|^2}I+\frac{i(w-\bar w)}{1+|w|^2}J+\frac{w+\bar w}{1+|w|^2}K$，
并补充 $I_0=I$、$I_\infty=-I$。每个 $I_w=aI+bJ+cK$ 满足 $a^2+b^2+c^2=1$，都是与同一度规相容的可积复结构。这一族复结构的参数空间称 twistor $\mathbb P^1$：原文 §4.1 记作 $\mathbb P^1_h$（$h$ 指 hyperkähler），本词典按术语表称「复结构 $\mathbb P^1_g$」；§5 中 $\mathbb P^1_g$ 又记四维 twist 参数 $t$ 的空间，靠嵌入 $w_+=-t$ 与之相接。}
\mathex{$G=SU(2)$：$I_w$ 使 $A_{\bar z}-w\phi_{\bar z}$ 与 $A_z+w^{-1}\phi_z$ 全纯。特别地 $w=0$ 给 $I$（Higgs 丛），$w=\infty$ 给 $-I$，$w=-i$ 给 $J$（平坦丛），$w=1$ 给 $K$。$U(1)$ 作用 $\phi\to\lambda\phi$（$|\lambda|=1$）把 $I_w$ 送到 $I_{\lambda^{-1}w}$，并固定 $w=0,\infty$ 两点。}
\physdef{物理上，$\mathbb P^1_h$ 的两份拷贝参数化二维 sigma 模型的拓扑场论：一对复结构 $(I_{w_+},I_{w_-})$ 决定一个 TQFT，B-模型要求 $w_+=w_-$，A-模型要求 $w_-=-1/\bar w_+$。四维理论给出的族由 twist 参数 $t\in\mathbb P^1_g$ 参数化，并嵌入为 $w_+=-t$、$w_-=t^{-1}$（式 (5.11)）。}
\physex{$t=\pm i$ 时 $w_+=w_-=\mp i$，得到复结构 $J$ 或 $-J$ 中的 B-模型；$t=\pm1$ 时得到复结构 $K$ 或 $-K$ 中的 A-模型；$t$ 为实数时（含 $t=0,\infty$）得到 A-模型。$t$ 与 $-t$ 给出二维等价的 TQFT，这来自 $SU(4)_{\mathcal R}$ 中心的作用。}
\bridge{twistor $\mathbb P^1$ 与四维 twist 参数空间的对应 $w_+=-t$、$w_-=t^{-1}$ 把「三个复结构」翻译成「$\mathbb P^1_g$ 上的点」：$t=\pm i$ 落在 B-模型支，实 $t$ 落在 A-模型支（§5.1--§5.2）。}
\intuit{把 twistor $\mathbb P^1$ 想成「复结构的罗盘」：$I,J,K$ 是三个基本方向，$I_w$ 是所有单位线性组合；物理上转罗盘就是选不同的扭化与超荷。}
\pitfall{它不是靶空间的一族度量，也不是一族不同的流形：底流形与度规固定，只有复结构在变。别把 $I_w$ 与 Penrose 的 twistor 空间（复 3 维、以 $\mathbb P^1$ 为纤维）混为一谈；也别把 $w$ 与物理 twist 参数 $t$ 直接等同——两者差一个分式嵌入 $w_+=-t$、$w_-=t^{-1}$。}
\end{word}

\begin{word}{辛形式}{symplectic form}{symplectic-form}{§4.1 p.44}
\mathdef{超凯勒流形上三个辛形式由度规给出：$\omega_I(\cdot,\cdot)=g(I\cdot,\cdot)$、$\omega_J(\cdot,\cdot)=g(J\cdot,\cdot)$、$\omega_K(\cdot,\cdot)=g(K\cdot,\cdot)$。它们都是闭的、非退化的 2-形式；$\Omega_I=\omega_J+i\omega_K$ 关于 $I$ 为 $(2,0)$ 型闭形式（全纯辛形式），其余由循环置换得到 $\Omega_J=\omega_K+i\omega_I$、$\Omega_K=\omega_I+i\omega_J$。反过来，$\Omega_I$ 的实部与虚部给出 $\omega_J$ 与 $\omega_K$。}
\mathex{本文式 (4.8) 给出 $\mathcal U$ 上的显式形式：$\omega_I=-\frac{1}{4\pi}\int_C\mathrm{Tr}(\delta A\wedge\delta A-\delta\phi\wedge\delta\phi)$，$\omega_J=\frac{1}{2\pi}\int|d^2z|\,\mathrm{Tr}(\delta\phi_{\bar z}\wedge\delta A_z+\delta\phi_z\wedge\delta A_{\bar z})$，$\omega_K=\frac{1}{2\pi}\int\mathrm{Tr}\,\delta\phi\wedge\delta A$。其中 $\omega_J,\omega_K$ 恰当（$\omega_J=\delta\lambda_J$），$\omega_I$ 不恰当，但 $[\omega_I/2\pi]$ 是整上同调类，生成 $H^2(\mathcal M_H,\mathbb Z)$（单连通 $G$）。}
\physdef{三个辛形式在物理上分别是：$\omega_I$ 给出二维 sigma 模型的 $B$-场 $B=-(\theta/2\pi)\omega_I=-(\mathrm{Re}\,\tau)\omega_I$（来自四维 $\theta$ 项）；$\omega_K$ 是 A-模型的辛形式（A-模型只依赖辛形式，不依赖复结构）；$\omega_K$ 与 $\omega_I$ 组成复结构 $J$ 的全纯辛形式 $\Omega_J=\omega_K+i\omega_I$，而 $\omega_J,\omega_K$ 组成结构 $I$ 的 $\Omega_I$。}
\physex{$G=SU(2)$、$C$ 亏格 2：$\omega_J,\omega_K$ 恰当说明对应的 $B$-场拓扑平凡、可被重新定义吸收；$\omega_I$ 的整上同调类则意味着 $B$-场有 $\theta$ 周期性，且 $\int_\Sigma\Phi^*(\omega_I/2\pi)$ 是整数，与四维瞬子数同源。}
\bridge{物理侧的 $\theta$ 项、$B$-场与 A-模型的辛结构都由数学侧的 $\omega_I,\omega_J,\omega_K$ 给出：$B\propto\omega_I$、A-模型用 $\omega_K$、复结构 $I$ 的全纯辛形式是 $\Omega_I=\omega_J+i\omega_K$（§4.1、§5）。}
\intuit{辛形式 = 「给每对切向量算一个面积的尺子」：三把尺子属于三副复结构，但在同一度规下兼容，所以能互相转换。直觉靠得住，是因为超凯勒条件要求三把尺子都闭。}
\pitfall{不要以为三个辛形式地位完全对称：$\omega_J,\omega_K$ 恰当而 $\omega_I$ 不恰当（其上同调类非零，对应行列式线丛 $c_1$）。也不要把 Kähler 形式 $\omega_I$ 与 $I$ 的全纯辛形式混淆——后者是 $\Omega_I=\omega_J+i\omega_K$；更不要把 A-模型用 $\omega_K$、$B$-场用 $\omega_I$ 的角色弄反。}
\end{word}

\begin{word}{moment map $\mu$}{moment map}{moment-map}{§4.1 p.45}
\mathdef{设 Lie 群 $H$ 以辛方式作用在辛流形 $(X,\omega)$ 上，$V(\epsilon)$ 是生成元 $\epsilon\in\mathfrak h$ 对应的向量场。若存在映射 $\mu:X\to\mathfrak h^*$，使对一切 $\epsilon$ 有 $d\langle\mu,\epsilon\rangle=\iota_{V(\epsilon)}\omega$（右端是把 $V(\epsilon)$ 缩进辛形式所得的 1-形式），则称 $\mu$ 为 moment map。对超凯勒流形上保持 $I,J,K$ 的作用，三个辛形式各给一个 moment map $\mu_I,\mu_J,\mu_K$，合成分量向量 $\vec\mu=(\mu_I,\mu_J,\mu_K)$。}
\mathex{本文中 $\mathcal U$ 上作用着 $H=\mathcal G=\mathrm{Maps}(C,G)$：$\epsilon\in\Omega^0(C,\mathrm{ad}E)$ 生成的向量场为 $\delta A=-D\epsilon$、$\delta\phi=[\epsilon,\phi]$。式 (4.11) 给出 $\mu_I=-\frac{1}{2\pi}\int_C\mathrm{Tr}\,\epsilon\,(F-\phi\wedge\phi)$、$\mu_J=-\frac{1}{2\pi}\int_C|d^2z|\,\mathrm{Tr}\,\epsilon\,(D_z\phi_{\bar z}+D_{\bar z}\phi_z)$、$\mu_K=-\frac{i}{2\pi}\int_C|d^2z|\,\mathrm{Tr}\,\epsilon\,(D_z\phi_{\bar z}-D_{\bar z}\phi_z)$。取遍检验函数 $\epsilon$，三个 moment map 为零恰好写出三条 Hitchin 方程。}
\physdef{moment map 在物理上就是 D-项方程。二维 gauged sigma 模型里，超凯勒靶空间 $X$ 上的规范群作用给出三个 D-项，其超空间形式正是三个 moment map；解 D-项条件并积掉有质量的规范多重态，低能靶空间就是 $X/\!/\!/H$。故「moment map 零点」=「D-项方程」=「规范场低能被迫取到的值」。}
\physex{本文情形 $H=\mathrm{Maps}(C,G)$ 是四维理论紧化到 $C$ 后残存的二维规范群。$\mu_J=0$ 与 $\mu_K=0$ 合起来给 $D_z\phi_{\bar z}=D_{\bar z}\phi_z=0$（即 $D\phi=0$、$D\star\phi=0$ 的复组合），$\mu_I=0$ 给 $F=\phi\wedge\phi$——这就是原文式 (4.1) 的三条方程。}
\bridge{一一对应：三个 moment map 同时为零 ⟺ Hitchin 方程；商掉 $\mathcal G$ ⟺ 模规范变换。物理叙述「二维 $(4,4)$ 理论取低能极限、积掉规范多重态」与数学叙述「超凯勒商 $\mathcal U/\!/\!/\mathcal G$」是同一件事（§4.1 式 (4.11) 之后）。}
\intuit{moment map 是「守恒量在流形上的实现」：无穷小对称（向量场）缩进辛形式得到闭 1-形式，它必须恰当，位势就是 moment map。三个 moment map 同时取零 = 三个守恒量都取「基态值」。}
\pitfall{只把一个 moment map 归零是普通 Kähler 商，得到的不是 Hitchin 模空间；必须三个一起为零。另一坑：当作用带非零中心荷时零点应改为常值（仿射商），本文取零；若照搬「零点」而不检查中心荷，商出来的空间会差一个环面。}
\end{word}

\begin{word}{Hitchin 纤维化}{Hitchin fibration}{hitchin-fibration}{§4.3 p.53}
\mathdef{Hitchin 纤维化是映射 $\pi:\mathcal M_H\to B$，$(E,\varphi)\mapsto$ $\varphi$ 的规范不变量多项式（特征多项式的系数）。具体地：$G=SU(2)$ 时 $\pi=\mathrm{Tr}\,\varphi^2\in H^0(C,K_C^2)$；$G=SU(N)$ 时取 $w_n=\mathrm{Tr}\,\varphi^n$（$n=2,\dots,N$）；一般 $G$ 取 $r=\mathrm{rank}\,G$ 个独立不变齐次多项式 $\mathcal P_i$（次数 $d_i$），$B=\bigoplus_{i=1}^r H^0(C,K_C^{d_i})$。基 $B$ 称 Hitchin 基，$\pi$ 使它成为代数完全可积系统。}
\mathex{$G=SU(2)$、$C$ 亏格 $g\ge2$：$B=H^0(C,K_C^2)$，$\dim_{\mathbb C}B=3g-3$；$g=2$ 时 $\dim B=3$ 而 $\dim_{\mathbb C}\mathcal M_H=6$，故一般纤维是复 3 维紧环面。$G=SU(3)$、$g=2$：$\dim B=3+5=8$，而 $\dim_{\mathbb C}\mathcal M_H=2\cdot 8=16$，纤维是复 8 维环面，恰好一半一半。$g=0$、$C=\mathbb P^1$、$G=SU(2)$ 且度为 $0$：$K_C^2$ 无整体截面，Hitchin 方程只剩 $\varphi=0$ 与平凡丛，$\mathcal M_H$ 退化成一点——这是最干净的边界情形。}
\physdef{在物理侧，$B$ 是把 Higgs 场「对角化」后剩下的规范不变量：$\varphi$ 的特征值（即 $C$ 上各点的质量/真空参数）。它不是额外自由度，而是低能理论里可同时测量的那半数自由度；另外一半藏在纤维（环面）上。}
\physex{镜像对称的 SYZ 图像：四维理论的两个对偶描述紧化到 $C$ 后，$\mathcal M_H(G)$ 与 $\mathcal M_H({}^LG)$ 纤维化到同一个 $B$（双 Hitchin 纤维化），一般纤维互为对偶环面；$B$-模型 on $\mathcal M_H({}^LG)$ 与 A-模型 on $\mathcal M_H(G)$ 的等价由纤维上的 T-对偶逐点实现（§5）。基 $B$ 正是这对镜像理论共享的参数空间。}
\bridge{数学侧：$\pi$ 把 Higgs 丛映到特征多项式，是代数完全可积系统的动量映射；物理侧：$B$ 是镜像两侧共同的对偶基，纤维上的 T-对偶给出 A/B 模型等价（§5）。$T^*\mathcal M\hookrightarrow\mathcal M_H$ 给出另一个方向的纤维化（见「Hitchin 模空间」条）。}
\intuit{Hitchin 纤维化像「把矩阵对角化」：$\varphi$ 的特征值（特征多项式）规范不变，构成基 $B$；给定特征值后剩下的自由度（特征向量沿 $C$ 的 twisting）构成纤维，一般是环面。}
\pitfall{最易混的是把 Hitchin 纤维化与「忘掉 $\varphi$ 的第二纤维化」$\mathcal M_H\to\mathcal M$ 弄混：后者靶空间是丛模空间 $\mathcal M$（$T^*\mathcal M$ 在其 Zariski 开集中），前者靶是 Casimir 多项式向量空间 $B$。另注意 $\dim B=\frac12\dim\mathcal M_H$，不要指望 $\pi$ 是双有理同构。}
\end{word}

\begin{word}{可积系统}{integrable system}{integrable-system}{§4.3 p.54}
\mathdef{复完全可积系统是三元组 $(X,\Omega,\pi)$：$X$ 为复辛流形（$\dim_{\mathbb C}X=2n$），$\Omega$ 为全纯辛形式，$\pi:X\to B$ 为全纯映射，$\dim_{\mathbb C}B=n$，且 $\pi$ 的纤维是 $\Omega$ 下的 Lagrange 子流形（切空间全迷向且维数最大）。等价说法：$B$ 上局部存在 $n$ 个独立、两两 Poisson 交换的全纯函数，它们给出 $n$ 个交换的哈密顿流，可同时积分。}
\mathex{Hitchin 1987b：$\mathcal M_H\to B$ 是复完全可积系统，$\dim_{\mathbb C}\mathcal M_H=2\dim_{\mathbb C}B$，一般纤维是紧复环面、在 $\Omega_I$ 下是 Lagrange 子流形（对 $\omega_J,\omega_K$ 亦然，§4.3）。$G=SU(2)$、$g=2$ 时 $\dim B=3$，纤维是复 3 维环面；沿基 $B$ 移动时角变量在环面上线性演化。另有 $H_a=\int_C\alpha_a\wedge\mathrm{Tr}\,\varphi^2$（$\alpha_a$ 取遍 $H^1(C,T_C)$ 的基，$a=1,\dots,3g-3$）给出的交换流。}
\physdef{物理上，可积性意味着四维扭化理论紧化到 $C$ 后得到的二维理论具有足够多的守恒量：基 $B$ 的坐标是一组对合的守恒量，理论可用作用量-角变量求解，且 Hitchin 纤维的演化是线性的。}
\physex{可积性也是镜像对称能「逐点」实现的前提：每根纤维是环面，而环面上的 T-对偶就是 Fourier--Mukai 变换——这正是把 $\mathcal M_H(G)$ 的 A-模型与 $\mathcal M_H({}^LG)$ 的 B-模型对起来的机制（§5）。故「可积」在物理上不是数学装饰，而是 SYZ 图像的几何基础。}
\bridge{$B$ 的维数恰为 $\mathcal M_H$ 的一半，且纤维是 Lagrange 环面，这就是「可积」的全部内容；物理侧的守恒量、作用量-角变量与数学侧的 Casimir 基、环面纤维逐项对应（§4.3）。}
\intuit{可积系统 = 「维数正好一半的环面纤维」：$\dim$ 纤维 $=\dim$ 基 $=\frac12\dim X$，且纤维在辛形式下全迷向。直觉的边界：可积只保证有足够多守恒量、能写成作用量-角变量，\textbf{不}保证能用初等函数写出解。}
\pitfall{「可积」不等于「可解」：显式解一般要用 theta 函数，奇异纤维处结构复杂（这正是几何 Langlands 中几何内窥的入口）。另一常见错误是把 $\mathcal M_H$ 本身当成相空间——在二维约化里它是靶空间/构型空间，$T^*\mathcal M_H$ 才更像相空间；混淆二者会把 Lagrange 条件用错对象。}
\end{word}

\begin{word}{Casimir 基（Hitchin 基 $B$）}{Casimir basis}{casimir-basis}{§4.3 p.53}
\mathdef{$B=\bigoplus_{i=1}^{r}H^0(C,K_C^{d_i})$，$r=\mathrm{rank}\,G$，$d_i$ 是 $G$ 的 Casimir 不变量的次数（$d_i=m_i+1$，$m_i$ 为指数），满足 $\sum_{i=1}^{r}(2d_i-1)=\dim G$。它是 $\varphi$ 的基不变量多项式的靶空间：$\pi(E,\varphi)=(\mathcal P_1(\varphi),\dots,\mathcal P_r(\varphi))$。更内蕴的写法是 $B=H^0(C,(\mathfrak g//G)\otimes K_C)$，其中 $\mathfrak g//G=\mathrm{Spec}((\mathrm{Sym}\,\mathfrak g^*)^G)$ 是 Casimir 环的谱。}
\mathex{$G=SU(2)$（$g\ge2$）：$r=1$，唯一基本不变量是 $\mathrm{Tr}\,\varphi^2$，$d=2$，$B=H^0(C,K_C^2)$，$\dim=3g-3$。$G=SU(3)$：$d_1=2,d_2=3$，$B=H^0(K_C^2)\oplus H^0(K_C^3)$，$\dim=3(g-1)+5(g-1)=8(g-1)$，而 $\dim_{\mathbb C}\mathcal M_H=(2g-2)\cdot 8=16(g-1)$，恰是一半。一般 $SU(N)$（$g\ge2$）：$d_i=i+1$（$i=1,\dots,N-1$），$\sum_{i=2}^{N}(2i-1)=N^2-1=\dim SU(N)$，$\dim B=(g-1)(N^2-1)$。}
\physdef{物理上 Casimir 不变量就是 Higgs 场的可观测规范不变量——把 $\varphi$ 在各点的特征值组合成不变多项式。$B$ 的坐标是二维有效理论中可同时测量的守恒量；$B$ 的维数只有 $\mathcal M_H$ 的一半，正是「可积」的物理含量。}
\physex{$G=SU(2)$、$g=2$：$B$ 是 3 维复向量空间，其坐标即三个二次微分的系数，物理上对应 $\mathrm{Tr}\,\varphi^2$ 的模；镜相对偶下 $SU(2)$ 自对偶，两侧的 $B$ 由不变双线性型自然等同。}
\bridge{数学侧：$B$ 是 Hitchin 纤维化的基，也是 Chevalley 定理 $\mathfrak g//G=\mathfrak h//W$ 的几何化身；物理侧：$B$ 是镜像两侧共享的参数空间。注意 ${}^LB$ 与 $B$ 的可典范等同依赖不变双线性型 $\kappa_0$ 的选择（用它识别 $\mathfrak h$ 与 $\mathfrak h^*$），换 $\kappa_0$ 只差一个 $\mathbb C^\times$ 作用，可由 $\varphi$ 的标度吸收。}
\intuit{Casimir 基 = 「拿特征多项式的系数当坐标」：对矩阵就是 $(\mathrm{tr}\,A,\mathrm{tr}\,A^2,\dots)$；在模空间上这些值把 Higgs 丛压成一张「特征值分布」的影像。直觉边界：$\mathcal P_i$ 只捕获特征值，不捕获特征向量——后者才是纤维的自由度。}
\pitfall{不要把 $B$ 与 $\mathcal M_H$ 的复化混同；也不要以为 $d_i$ 可以任取：它由 $G$ 的指数决定并满足 $\sum(2d_i-1)=\dim G$，于是 $\dim B=(g-1)\dim G=\frac12\dim\mathcal M_H$ 是恒等式而非巧合。另一坑：把 $B$ 与 ${}^LB$ 的等同当成绝对规范——它依赖 $\kappa_0$ 的选择。}
\end{word}

\begin{word}{谱覆盖}{spectral cover}{spectral-cover}{§4.3 p.54}
\mathdef{设 $G=SU(N)$，$C$ 为亏格 $g\ge2$ 的紧黎曼面，$E\to C$ 为 $SU(N)$-丛，$\varphi\in H^0(C,K_C\otimes\mathrm{End}E)$。谱覆盖是 $B$ 上的曲线族 $\Sigma_b\subset K_C$，由特征多项式定义：$\Sigma_b=\{\lambda\in K_C:\ \det(\lambda-\varphi)=0\}$，展开为 $\lambda^N+s_2\lambda^{N-2}+\cdots+s_N=0$（$s_n$ 是 $\mathrm{Tr}\,\varphi^n$ 的对称化）。投影 $\Sigma_b\to C$ 是 $N$ 重覆盖，分支除子由判别式给出；一般 $b$ 时 $\Sigma_b$ 光滑，亏格 $g_\Sigma=N^2(g-1)+1$（$N=2$ 时 $g_\Sigma=4g-3$）。}
\mathex{$G=SU(2)$、$g=2$：$\Sigma_b$ 是 $K_C$ 中亏格 5 的二重覆盖，Hitchin 纤维一般同构于 $\mathrm{Prym}(\Sigma_b/C)$（或它的 torsor），$\dim\mathrm{Prym}=g_\Sigma-g=3g-3=3$，与 $\dim B$ 相同。一般 $N$：纤维是 $\Sigma_b$ 的广义 Prym 簇，维数 $(N^2-1)(g-1)=\dim B$。这就是「特征值 + 谱曲线上的线丛 = Higgs 丛」的精确含义。}
\physdef{谱覆盖在物理侧是把 Hitchin 基 $B$ 的坐标「几何化」：每个 $b$ 处，纤维环面被实现为谱曲线上的 Jacobi/Prym 簇；镜像对称的 T-对偶就作用在这些环面上，把 $\mathcal M_H(G)$ 的 A-模型与 $\mathcal M_H({}^LG)$ 的 B-模型逐点对起来。}
\physex{$G=SU(2)$、$g=2$：谱覆盖是亏格 5 的二重覆盖，Hitchin 纤维是 3 维 Prym 环面；$SU(2)$ 自对偶，镜像两侧的纤维互为对偶环面，对偶在谱覆盖上由线丛的平移实现。}
\bridge{数学侧：谱覆盖给出 Hitchin 纤维的具体模型（Jacobi/Prym 簇），是可积系统显式化的标准技术；物理侧：它是 SYZ 镜像中纤维环面的几何底座。但 KW06 明确声明本文不用这套构造（§4.3 p.54：「在本文中，我们将在不使用谱覆盖构造的前提下尽量推进」），故谱覆盖在本词典中是背景知识而非原文论证工具。}
\intuit{谱覆盖 = 「把 $\varphi$ 的特征值当成 $C$ 上的一张多层覆盖」：每点 $z\in C$ 上方 $\varphi(z)$ 有 $N$ 个特征值，于是 $C$ 上长出一张 $N$ 层覆盖 $\Sigma_b$；Hitchin 纤维是其上的线丛。直觉边界：谱覆盖本身是 $K_C$ 中的曲线，不是 $\mathcal M_H$ 的纤维。}
\pitfall{最易错的是把谱覆盖（$K_C$ 中的曲线 $\Sigma_b$）与 Hitchin 纤维（$\mathcal M_H$ 的纤维，复环面）混为一谈——纤维是 $\Sigma_b$ 的 Jacobi/Prym 簇，是环面而非曲线。另一坑：把「谱曲线」与「谱覆盖」当成两个不同对象——谱曲线即总空间 $\Sigma_b$，「覆盖」强调投影 $\Sigma_b\to C$。最后：原文声明不用此技术，别把谱覆盖的论证当成 KW06 的证明。}
\end{word}

\par\medskip
\noindent{\small 本轨道小结：Hitchin 模空间 $\mathcal M_H$ 是 4 维扭化 $\mathcal N=4$ 理论紧化到黎曼面 $C$ 的低能靶空间，同时是 Higgs 丛与平坦丛两套描述的同一实流形（超凯勒）；三个复结构 $I,J,K$ 与三个辛形式 $\omega_I,\omega_J,\omega_K$ 把 A/B 模型、$B$-场与 $\theta$ 项一一对应；Hitchin 纤维化 $\pi:\mathcal M_H\to B$ 把它变成完全可积系统，一般纤维是复环面，其上的 T-对偶正是镜像对称的几何机制，也是几何 Langlands 的出发点。}

